\note{1}{Thu 28 Dec 2023 21:53}{Brydges Exercise Lecture 1}

\begin{problem}
	Problem 1.1. An $n \times n$ matrix $A$ is said to be positive-definite if for all non-zero $\lambda=$ $\left(\lambda_1, \ldots, \lambda_n\right)$ in $\mathbb{C}^n$
\[
\sum_{1 \leq i, j \leq n} \lambda_i A_{i j} \bar{\lambda}_j>0
\]

If the inequality is not strict then the matrix is said to be positive-semidefinite. A continuous function $f: \mathbb{R} \rightarrow \mathbb{R}$ is said to be positive-definite if for every $x \in \mathbb{R}^n$, the matrix ( $f\left(x_i-\right.$ $\left.\left.x_j\right)\right)_{1 \leq i, j \leq n}$ is positive-semidefinite. Prove that:
(1) If $f$ is positive-definite, then $V$ given by a the two-body potential $v(x, y)=f(x-y)$ as in Example 1.2 satisfies the stability bound (1.5).
(2) If $f$ is continuous and integrable so that the Fourier transform $\hat{f}(k)=\int f(x) e^{-i k x} d x$ exists, if $\hat{f} \geq 0$, then $f$ is positive-definite. (This is the "easy" half of Bochner's theorem.) Concentrate on the case where $\hat{f}$ is also integrable, and then see exercise 8.4.30 in [Fol99] to remove this assumption.

This extends to $\mathbb{R}^d$. Conjecture: Every two-body potential $v(x, y)=f(x-y)$ such that $V$ satisfies (1.5) has the form $f=$ non-negative function + positive-definite function.
\end{problem}
\begin{proof}
	1) Consider  $1$-dimensional matrix, we have $f(0) \ge 0$. Now consider $2$-dimensional matrix, set $\lambda_1 = \lambda_2 = 1$, we have  $2 f(0) + f(x_1 - x_2) + f(x_2 - x_1) \ge 0$. Since the matrix $f(x_i - x_j)$ is positive semi-definite, it is Hermitian, so in particular symmetric. Hence $f(0) + f(x_1 - x_2) \ge 0$. $x$ is arbitrary, so $f \ge  - f\left( 0 \right) $.

	2) Assume $\hat{f} $ is integrable, then we can write $f(x_i - x_j)$ using inverse fourier transform:
	\[
		\sum_{i j} \lambda_i f(x_i - x_j) \overline{\lambda} _j = \int \hat{f} (k) \sum_{i,j} \lambda_i e^{i k (x_i - x_j)} \overline{\lambda} _j dk
	.\] 
	Now, note that for all $k$, $x \mapsto e^{i k x}$ is positive semidefinite, since
	\[
		\sum_{i j} \lambda_i e^{i k (x_i - x_j)} \overline{\lambda} _j
		= \sum_i \lambda_i e^{i k x_i} \sum_j \overline{\lambda} _j e^{- i k x_j}
		= \left| \sum_{i} \lambda_i e^{i k x_i} \right| ^2 \ge 0
	.\] 
	By assumption $\hat{f} (k) \ge 0$. Thus $f$ is positive semidefinite.

	If we strengthen the assumption to $\hat{f} > 0$ locally, then $f$ is indeed positive definite.
\end{proof}

\begin{definition}[Gibbs Measure]
	Definition 1.3. Let $\Lambda \subset \mathbb{R}^d,|\Lambda|<\infty$. Let $z \geq 0$. For $E \subset \Lambda^*$
\[
\mathbb{P}(E)=\frac{1}{Z} \sum_{n \geq 0} \frac{z^n}{n !} \int_{E \cap \Lambda^n} e^{-V} d x_1 \cdots d x_n
\]
is called the Grand Canonical Ensemble Gibbs measure, or simply the Gibbs measure, where
\[
Z=\sum_{n \geq 0} \frac{z^n}{n !} \int_{\Lambda^n} e^{-V} d x_1 \cdots d x_n
\]
is the normalization factor such that $\mathbb{P}\left(\Lambda^*\right)=1$.
\end{definition}

This can be understood as
\[
	d P(x) = \frac{1}{Z} \frac{z^n}{n!} e^{-V(x)} d x \qquad n = N(x)
.\] 

To take expectations, for $F: \Lambda^* \to \mathbb{R}$,
\[
	E F = \int _{\Lambda^*} F d P = \frac{1}{Z} \sum_{n \ge 0} \frac{z^n}{n!} \int_{\Lambda^n} e^{-V} F dx
.\] 

\begin{claim}
	If $V(x) = - \sum_{i = 1}^{N(x)} \phi(x_i)$, then
	\[
		Z = \exp\left( z \int  _\Lambda e^{\phi} dx \right) 
	.\] 
\end{claim}
\begin{proof}
	\begin{align*}
		Z&=\sum_{n \geq 0} \frac{z^n}{n !} \int_{\Lambda^n} \exp\left( \sum_{i = 1}^{n} \phi(x_i) \right)  d x_1 \cdots d x_n\\
		&=\sum_{n \geq 0} \frac{z^n}{n !}\left( \int_{\Lambda} \exp\left(\phi(x) \right)  d x\right)^n\\
		&= \exp\left(z \int _{\Lambda} e^{\phi} dx\right)
	.\end{align*}
\end{proof}





\begin{Lemma}[Ideal Gas]
	Let $V=0$. Let $X_1, \ldots, X_n \subset \Lambda$, where $|\Lambda|<\infty$. Then
a) $N\left(X_i\right) \sim \operatorname{Poisson}\left(z\left|X_i\right|\right)$;
b) if $\left|X_i \cap X_j\right|=0$ for $i \neq j$, then $N\left(X_1\right), \ldots, N\left(X_n\right)$ are independent.
\end{Lemma}
\begin{proof}
	We show that the generating function of $N(X_i)$ matches the generating function of the respecting Poisson distribution. The generating function of $Y \sim \text{Poisson}(r)$ is:
	\[
		E e^{t Y} 
		= \sum_{n \ge 0} \frac{r^n}{n!}e^{-r} e^{tn}
		= \sum_{n \ge 0} \frac{(re^t)^n}{n!}e^{-r}
		= e^{-r} e^{r e^{t} } 
	.\] 

	Now, the number of particles in path $x$ hitting region $X \subset \Lambda$ is
	\[
		N(X, x) = \sum_{i = 1}^{N(x)} \mathbf{1}(x_i \in X)
	.\] 
	The mgf is
	\[
		E(\exp\left(tN(X, \cdot ) \right) ) = \frac{1}{Z} \sum_{n \ge 0} \frac{z^n}{n!} \int _{\Lambda^n} \exp\left(t \sum_{i = 1}^n \mathbf{1}(x_i \in X)  \right) xdx = Z(\phi) / Z, \qquad \phi(x_i) = t \mathbf{1}(x_i \in X)
	.\] 
	By the preceding claim,
	\[
		E\left( \exp\left(t N(X) \right)  \right) 
		= \exp\left(z \int _{\Lambda} e^{\phi} dx \right) / \exp\left( z |\Lambda| \right) 
		= 
		\exp\left(z \int _{\Lambda} e^{\phi}  - 1 \right) 
		= \exp\left( z |X|(e^t - 1) \right) 
	.\] 
	This implies $ N(X) \sim \text{Poisson}(z |X|)$.

	For independence, we also use the fact that $X_1, \ldots, X_m$ independent iff their joint mgf splits.
	Note that
	\[
		E\left( \exp\left(t_1 N(X_1) + \ldots + t_m N(X_m) \right)  \right) 
		= Z(\phi) / Z
	\] 
	where $\phi(x_j) = \sum_{i = 1}^m t_i \mathbf{1}\left( x_j \in X_i \right) $.
	Hence, this is equal to
	\[
		\exp\left(z \int _{\Lambda} e^{\phi}  - 1 \right) 
	.\] 
	\todo{THe normalization is still not right.}
	
\end{proof}

\begin{problem}
	Find $\left\langle N(\Lambda) \right\rangle  / |\Lambda|$.
\end{problem}
\begin{proof}
	\begin{align*}
		\left\langle N(\Lambda) \right\rangle 
		&= \frac{1}{Z} \sum_{n \ge 0} \frac{z^n}{n!}\int _{\Lambda^n} n dx \\
		&= \frac{1}{Z} z |\Lambda| \exp\left(z |\Lambda| \right) 
	.\end{align*}
	But $Z = \exp\left(z \left| \Lambda \right|  \right) $. Hence
	\[
		\left\langle N(\Lambda) \right\rangle / \left| \Lambda \right|  = z
	.\] 
\end{proof}

Connection between Ising model and mean field theory with two states: Subcritical Ising gives same long-term behavior, as it locks into a universal state.

Connection between De Finetti Theorem: Distribution of infinite exchangeable sequence, conditioned on the exchangeable field, are iid Bernoulli, the parameter given by some distribution.
In mean field theory, the blocks are exchangeable since there's no local interactions and the potential is spatially symmetric. Hence, the state of each block (have particle or not) is a Bernoulli event, with parameter given by the two-value distribution $\rho$, consistent with our model.

